%!TEX program = xelatex

\documentclass[11pt,a4paper]{extarticle}

% ---------------------------------------------------------------------
%  LIBRERÍA Y FUENTES DESDE GITHUB (siempre la última versión)
%  https://github.com/nataliagamezbarea/apuntes-latex
%
%  Al compilar: descarga estilo.tex + fonts/ de GitHub a una carpeta
%  temporal, los usa, y al terminar la BORRA (no queda nada).
%  Necesita -shell-escape (ya está activado en la config de VS Code).
% ---------------------------------------------------------------------
\immediate\write18{rm -rf "/tmp/apuntes-latex-lib-\jobname" && mkdir -p "/tmp/apuntes-latex-lib-\jobname" && curl -fsSL https://codeload.github.com/nataliagamezbarea/apuntes-latex/tar.gz/refs/heads/main | tar -xz -C "/tmp/apuntes-latex-lib-\jobname" --strip-components=1}
\AtEndDocument{\immediate\write18{(sleep 3; rm -rf "/tmp/apuntes-latex-lib-\jobname") >/dev/null 2>&1 &}}

% Carpeta temporal con estilo.tex y fonts/ (se borra al terminar)
\newcommand{\bibliotecasDir}{/tmp/apuntes-latex-lib-\jobname/bibliotecas_latex}

% La librería de estilos (va ANTES de \begin{document})
\input{\bibliotecasDir/estilo}

\usepackage[spanish]{babel}

\begin{document}

\paletaOceano
\titulo{TEMA 02}{Teoria de la Gestalt}
\usarfondo

% ===================== 1 =====================
\secciongrande{Percebre no és rebre}

\parrafo{Segons la psicologia de la Gestalt, la ment configura mitjançant
certes lleis els elements que li arriben pels sentits (percepció) o per la
memòria (pensament).}

\lista{\textbf{Buscar}|\textbf{Seleccionar}|\textbf{Relacionar}|
       \textbf{Establir connexions}|\textbf{Recordar}|
       \textbf{Identificar}|\textbf{Jerarquitzar}|\textbf{Avaluar}|
       \textbf{Aprendre}|\textbf{Interpretar}}

% ===================== 2 =====================
\secciongrande{El missatge visual}

\parrafo{Dissenyar és construir un missatge amb una intenció concreta.}

\lista{\textbf{Buscar}|\textbf{Seleccionar}|\textbf{Relacionar}|
       \textbf{Establir connexions}|\textbf{Recordar}|
       \textbf{Identificar}|\textbf{Jerarquitzar}|\textbf{Avaluar}|
       \textbf{Aprendre}|\textbf{Interpretar}}

% ---- Pie de página ----
\pagefoot{Gestalt -- Tema 02}

\end{document}
